% CONS_3D: Αναλυτικό Θεωρητικό Υπόμνημα & Οδηγός Χρήσης
% Μεταγλώττιση με LuaLaTeX:  lualatex cons3d_manual.tex  (δύο φορές)
% Το PDF αντιγράφεται στη ρίζα του αποθετηρίου ως «Εγχειρίδιο CONS3D.pdf».
\documentclass[11pt,a4paper]{article}

\usepackage{fontspec}
\setmainfont{Noto Sans}
\usepackage{polyglossia}
\setmainlanguage{greek}
\setotherlanguage{english}

\usepackage[a4paper,margin=2cm,top=2.5cm,bottom=2.5cm]{geometry}
\usepackage{amsmath}
\usepackage{booktabs}
\usepackage{enumitem}
\usepackage{tikz}
\usepackage{caption}
\usepackage{float}
\setlist[itemize]{label=-}
\captionsetup{labelsep=colon}
\setlength{\parindent}{0pt}
\setlength{\parskip}{6pt}
\setlength{\emergencystretch}{2em}

\newcommand{\cons}{\textbf{CONS\_3D}}

\title{\textbf{CONS\_3D: Αναλυτικό Θεωρητικό Υπόμνημα \& Οδηγός Χρήσης}}
\author{Ανάπτυξη: Π. Πελέκης, Αναπλ. Καθηγητής Π.Π.}
\date{Έκδοση: Νοέμβριος 2026}

\begin{document}
\maketitle

Το \cons{} είναι ένα προηγμένο εργαλείο αριθμητικής ανάλυσης για τον υπολογισμό των τάσεων, των καθιζήσεων και της χρονικής εξέλιξης της στερεοποίησης σε τρισδιάστατα γεωτεχνικά προβλήματα.

% =====================================================================
\section{Μέρος 1: Θεωρητικό Υπόβαθρο \& Αριθμητικές Μέθοδοι}

\subsection{Χωρική Διακριτοποίηση \& Γεωμετρία Μοντέλου}

Η επίλυση στο \cons{} βασίζεται στη διακριτοποίηση του συνεχούς εδαφικού μέσου σε έναν τρισδιάστατο κάνναβο (3D Grid) με διαστάσεις $N_x \times N_y \times N_z$ και βήματα $\Delta x, \Delta y, \Delta z$. Για την επίλυση της στερεοποίησης ο κάνναβος έχει $31 \times 41 \times 31$ κόμβους, δηλαδή πυκνότερη διακριτοποίηση κατά το βάθος ($y$), όπου οι υπερπιέσεις μεταβάλλονται πιο απότομα.

Η γεωμετρία της στρωματογραφίας στο επίπεδο X-Y ορίζεται ελεύθερα από τον χρήστη μέσω \textbf{κλειστών πολυγώνων (Polygons)}. Αυτή η προσέγγιση είναι εξαιρετικά ισχυρή, καθώς επιτρέπει την προσομοίωση \textbf{κεκλιμένων στρώσεων, φακών, ανισοπαχών αποθέσεων ή κεκλιμένου βραχώδους υποβάθρου}, σε αντίθεση με τις κλασικές μεθόδους που περιορίζονται σε αυστηρά οριζόντια επίπεδα. Τα όρια του καννάβου ορίζονται από τις στρώσεις· τα 2D επιχώματα (φορτία πάνω από την επιφάνεια) και οι γραμμές στράγγισης δεν τα επεκτείνουν.

\subsubsection{Διακριτοποίηση Επιφανειακών Φορτίων}

Για τον υπολογισμό των πρόσθετων τάσεων από σύνθετα σχήματα (ορθογώνια, κύκλους, τραπεζοειδή επιχώματα), το πρόγραμμα εκτελεί \textbf{αριθμητική ολοκλήρωση (Numerical Integration)}. Η φορτιζόμενη επιφάνεια τεμαχίζεται σε στοιχειώδη εμβαδά $dA = \Delta x \cdot \Delta z$, με ομοιόμορφη υποδιαίρεση που καλύπτει ακριβώς το αποτύπωμα κάθε φορτίου. Στο κυκλικό φορτίο τα στοιχειώδη φορτία κανονικοποιούνται, ώστε το συνολικό φορτίο να είναι ακριβώς $q\,\pi R^2$.

Κάθε στοιχειώδες εμβαδόν θεωρείται ότι ασκεί ένα συγκεντρωμένο σημειακό φορτίο:
\begin{equation}
  P_i = q(x,z)\cdot dA
\end{equation}

Η συνολική πρόσθετη τάση $\Delta\sigma_v$ σε ένα βάθος $y$ προκύπτει από την επαλληλία (superposition) των λύσεων του \textbf{Boussinesq} για όλα τα $P_i$:
\begin{equation}
  \Delta\sigma_v(x,y,z) = \sum_{i=1}^{N_{loads}} \frac{3P_i}{2\pi y^2}\left[\frac{1}{1+(r_i/y)^2}\right]^{5/2}
\end{equation}
Όπου $r_i$ η οριζόντια απόσταση μεταξύ του σημείου ενδιαφέροντος και του σημειακού φορτίου $P_i$.

Τα 2D επιχώματα, που εκτείνονται σε όλο το μήκος Z, τεμαχίζονται σε κατακόρυφες λωρίδες πλάτους $b$, με φορτίο $q = \gamma h$ από το ύψος $h$ στο μέσο κάθε λωρίδας. Η τάση από κάθε λωρίδα, με άκρα σε οριζόντιες αποστάσεις $x_1, x_2$ από το σημείο και σε βάθος $y$, δίνεται από την αναλυτική λύση του λωριδωτού φορτίου:
\begin{equation}
  \Delta\sigma_v = \frac{q}{\pi}\left[\arctan\frac{x_2}{y} - \arctan\frac{x_1}{y} + \frac{x_2\,y}{x_2^2+y^2} - \frac{x_1\,y}{x_1^2+y^2}\right]
\end{equation}

\subsubsection{Είδη Επιφανειακών Φορτίων \& Γεωμετρία}

Το πρόγραμμα υποστηρίζει πολλαπλούς τύπους επιφανειακών φορτίων, προσφέροντας πλήρη ευελιξία στην προσομοίωση πραγματικών κατασκευών:
\begin{itemize}
  \item \textbf{Ορθογωνικά \& Κυκλικά Φορτία:} Αντιπροσωπεύουν συνήθεις θεμελιώσεις, πέδιλα, κοιτοστρώσεις ή δεξαμενές. Για τα φορτία αυτά, ο χρήστης εισάγει απευθείας την τιμή της επιβαλλόμενης τάσης $q$ (σε kPa) και το εικονικό ειδικό βάρος $\gamma$ (σε kN/m$^3$). Το πρόγραμμα υπολογίζει \textbf{έμμεσα} το φυσικό ύψος του φορτίου ($h = q/\gamma$), το οποίο και χρησιμοποιείται αποκλειστικά για την ρεαλιστική 3D οπτική αναπαράσταση της κατασκευής πάνω στο εδαφικό ανάγλυφο.
  \item \textbf{Επιχώματα (2D \& 3D):} Προσομοιώνουν γραμμικά έργα (π.χ. οδοποιία, αναχώματα) ή τραπεζοειδείς αποθέσεις. Εδώ, η γεωμετρία ελέγχεται πλήρως από τον χρήστη: ορίζεται απευθείας το μέγιστο \textbf{ύψος} του επιχώματος καθώς και η \textbf{κλίση των πρανών} (εκφρασμένη σε λόγο Κατακόρυφο προς Οριζόντιο, V:1H). Ο αλγόριθμος αναλαμβάνει να προσαρμόσει αυτόματα το πλάτος στέψης και να υπολογίσει το ακριβές, μεταβαλλόμενο φορτίο $q(x,z)$ σε κάθε σημείο της βάσης του, προτού προχωρήσει στην αριθμητική ολοκλήρωση.
\end{itemize}

\subsection{Υπολογισμός Καθιζήσεων από Συμπιέσεις}

Η καθίζηση δεν υπολογίζεται με μία ενιαία εμπειρική σχέση, αλλά προκύπτει από την \textbf{κατακόρυφη ολοκλήρωση των παραμορφώσεων} κατά μήκος του εδαφικού προφίλ.

Για μια κατακόρυφο στη θέση $(x,z)$, το εδαφικό προφίλ διαιρείται σε $N_y$ στοιχειώδεις στρώσεις πάχους $\Delta y$.

\textbf{Βήμα 1ο: Υπολογισμός Αρχικής Κατάστασης}\\
Σε κάθε βάθος $y$, υπολογίζεται η αρχική ενεργός γεωστατική τάση:
\begin{equation}
  \sigma'_{v0} = \sum(\gamma\cdot\Delta y) - u_0
\end{equation}
Η τάση προστερεοποίησης καθορίζεται από τον λόγο OCR: $p'_c = \sigma'_{v0}\cdot OCR$.

\textbf{Βήμα 2ο: Υπολογισμός μεταβολής δείκτη πόρων $\Delta e$}\\
Βάσει της πρόσθετης τάσης $\Delta\sigma_v$, η τελική τάση είναι $\sigma'_{vf} = \sigma'_{v0} + \Delta\sigma_v$. Η μεταβολή του δείκτη πόρων $\Delta e$ υπολογίζεται σε τρεις διακριτές περιπτώσεις:
\begin{itemize}
  \item \textbf{α) Επαναφόρτιση (Reloading) -- $\sigma'_{vf} \le p'_c$:} Η τελική τάση δεν υπερβαίνει την τάση προστερεοποίησης και ο υπολογισμός γίνεται με τον δείκτη $C_s$:
  \begin{equation}\label{eq:reload}
    \Delta e = C_s \log\left(\frac{\sigma'_{vf}}{\sigma'_{v0}}\right)
  \end{equation}
  \item \textbf{β) Παρθενική Φόρτιση (Virgin Loading) -- $p'_c \le \sigma'_{v0}$:} Το έδαφος είναι πρακτικά κανονικά στερεοποιημένο (NC) και ο υπολογισμός γίνεται εξ ολοκλήρου με τον δείκτη συμπίεσης $C_c$:
  \begin{equation}
    \Delta e = C_c \log\left(\frac{\sigma'_{vf}}{\sigma'_{v0}}\right)
  \end{equation}
  \item \textbf{γ) Ενδιάμεση Κατάσταση -- $\sigma'_{v0} < p'_c < \sigma'_{vf}$:} Η φόρτιση ξεκινά από τον κλάδο επαναφόρτισης και, ξεπερνώντας το $p'_c$, μεταβαίνει στην παρθενική φόρτιση:
  \begin{equation}\label{eq:mixed}
    \Delta e = C_s \log\left(\frac{p'_c}{\sigma'_{v0}}\right) + C_c \log\left(\frac{\sigma'_{vf}}{p'_c}\right)
  \end{equation}
\end{itemize}

\textbf{Βήμα 3ο: Υπολογισμός Καθίζησης Στρώσης \& Ολοκλήρωση}\\
Η παραμόρφωση $\epsilon_y$ κάθε στοιχειώδους στρώσης είναι:
\begin{equation}
  \Delta S_i = \epsilon_{y,i}\cdot\Delta y = \frac{\Delta e_i}{1+e_{0,i}}\,\Delta y
\end{equation}
Η συνολική επιφανειακή καθίζηση $S_{surf}$ είναι το άθροισμα:
\begin{equation}
  S_{surf} = \sum_{i=1}^{N_y} \Delta S_i
\end{equation}

\subsubsection{Θεωρία Χαρακτηριστικού Σημείου (Kany) για Άκαμπτες Θεμελιώσεις}\label{sec:kany}

Όταν ένα θεμέλιο είναι \textbf{απολύτως εύκαμπτο} (π.χ. ένα επίχωμα ή μια δεξαμενή χωρίς άκαμπτο πυθμένα), η τάση επαφής με το έδαφος είναι ομοιόμορφη ($q =$ σταθ.), αλλά η καθίζηση παίρνει τη μορφή ``σκάφης'' (μεγαλύτερη στο κέντρο, μικρότερη στα άκρα). Αντίθετα, ένα \textbf{άκαμπτο θεμέλιο} (π.χ. κοιτόστρωση από οπλισμένο σκυρόδεμα) επιβάλλει \textbf{ομοιόμορφη καθίζηση} ($S =$ σταθ.). Για να συμβεί αυτό σε ένα ελαστικό μέσο, η τάση επαφής μεταβάλλεται ραγδαία, τείνοντας θεωρητικά στο άπειρο στα άκρα (λόγω συγκέντρωσης τάσεων) και ούσα ελάχιστη στο κέντρο:
\begin{equation}
  q(x) = \frac{2\,q_{avg}}{\pi\sqrt{1-\left(\frac{2x}{B}\right)^2}} \qquad \text{(για λωριδωτό θεμέλιο)}
\end{equation}

Ο αυστηρός υπολογισμός της καθίζησης ενός άκαμπτου πεδίλου απαιτεί την επίλυση ενός πολύπλοκου μητρώου αλληλεπίδρασης εδάφους-κατασκευής. Ωστόσο, ο \textbf{Kany (1959)} και άλλοι ερευνητές απέδειξαν μια εξαιρετικά κομψή ιδιότητα:
\begin{quote}
  Υπάρχει ένα συγκεκριμένο σημείο στην επιφάνεια του εδάφους όπου η καθίζηση ενός \textbf{εύκαμπτου} θεμελίου (υπό ομοιόμορφη φόρτιση $q_{avg}$) ισούται ακριβώς με την ομοιόμορφη καθίζηση ενός \textbf{άκαμπτου} θεμελίου των ίδιων διαστάσεων και του ίδιου συνολικού φορτίου. Αυτό ονομάζεται \textbf{Χαρακτηριστικό Σημείο (Characteristic Point)}.
\end{quote}

Η μαθηματική απόδειξη προκύπτει από την εξίσωση των ολοκληρωμάτων των δύο περιπτώσεων. Για ένα ορθογωνικό πέδιλο διαστάσεων $B\times L$, η αναλυτική λύση της θεωρίας ελαστικότητας δείχνει ότι το σημείο αυτό βρίσκεται στις συντεταγμένες:
\begin{equation}
  x_c \approx 0.74\,\frac{B}{2} \quad \text{και} \quad z_c \approx 0.74\,\frac{L}{2}
\end{equation}
Αντίστοιχα, για κυκλικό πέδιλο ακτίνας $R$, το χαρακτηριστικό σημείο βρίσκεται σε ακτινική απόσταση:
\begin{equation}
  r_c \approx 0.848\,R
\end{equation}

Στο \cons, η επιλογή ``Άκαμπτη Θεμελίωση'' εκμεταλλεύεται αυτό ακριβώς το θεώρημα. Το πρόγραμμα υπολογίζει την καθίζηση θεωρώντας εύκαμπτο φορτίο \textbf{μόνο στις συντεταγμένες του χαρακτηριστικού σημείου}, εξοικονομώντας τεράστιο υπολογιστικό χρόνο, και εν συνεχεία επιβάλλει αυτή την τιμή ομοιόμορφα σε όλο τον κάνναβο κάτω από το θεμέλιο. Η διόρθωση αυτή αφορά μόνο το μέγεθος της καθίζησης και δεν επηρεάζει τον ρυθμό στερεοποίησης (§\ref{sec:cv}).

\subsubsection{Τρισδιάστατη Επίδραση \& Συντελεστής Skempton-Bjerrum ($\mu$)}\label{sec:mu}

Η κλασική μονοδιάστατη θεωρία στερεοποίησης του Terzaghi υποθέτει ότι η αρχική υπερπίεση πόρων $\Delta u_0$ ισούται ακριβώς με την αύξηση της κατακόρυφης τάσης $\Delta\sigma_v$. Ωστόσο, σε πραγματικά τρισδιάστατα προβλήματα, η πλευρική παραμόρφωση δεν είναι μηδενική. Ο Skempton (1954) απέδειξε ότι η αναπτυσσόμενη υπερπίεση εξαρτάται από το πεδίο των τάσεων και την παράμετρο $A$ του εδάφους:
\begin{equation}
  \Delta u_0 = \Delta\sigma_3 + A(\Delta\sigma_1 - \Delta\sigma_3)
\end{equation}

Για να ληφθεί υπόψη αυτό το φαινόμενο, οι Skempton \& Bjerrum (1957) εισήγαγαν τον διορθωτικό συντελεστή $\mu$, ο οποίος συνδέει την πραγματική (τρισδιάστατη) καθίζηση στερεοποίησης $S_{3D}$ με την υπολογιζόμενη μονοδιάστατη $S_{1D}$:
\begin{equation}
  S_{3D} = \mu\cdot S_{1D}
\end{equation}

Η τιμή του συντελεστή $\mu$ εξαρτάται άμεσα από την παράμετρο $A$ του Skempton (άρα και από τον λόγο υπερστερεοποίησης OCR) και από τη γεωμετρία του φορτίου. Στον Πίνακα~\ref{tab:mu} δίνονται τυπικές τιμές του $\mu$ ανάλογα με την ιστορία φόρτισης του εδάφους. Ο χρήστης εισάγει τον $\mu$ ως ιδιότητα κάθε στρώσης και το πρόγραμμα πολλαπλασιάζει με αυτόν τις καθιζήσεις των στοιχειωδών στρώσεων. Ο $\mu$ διορθώνει μόνο το \textbf{μέγεθος} της καθίζησης: ο συντελεστής $m_v$ της επίλυσης στερεοποίησης (§\ref{sec:cv}) υπολογίζεται χωρίς αυτόν, ώστε να μην αλλοιώνεται ο ρυθμός στερεοποίησης.

\begin{table}[H]
  \centering
  \caption{Τυπικές τιμές παραμέτρου $A$ και συντελεστή $\mu$ (κατά Skempton \& Bjerrum)}\label{tab:mu}
  \begin{tabular}{lccc}
    \toprule
    \textbf{Τύπος Αργίλου} & \textbf{OCR} & \textbf{Παράμετρος $A$} & \textbf{Συντελεστής $\mu$} \\
    \midrule
    Πολύ ευαίσθητες άργιλοι (Sensitive) & 1.0 & $1.0 \sim 1.5$ & $1.0 \sim 1.2$ \\
    Κανονικά στερεοποιημένες (NC) & 1.0 & $0.5 \sim 1.0$ & $0.7 \sim 1.0$ \\
    Ελαφρώς υπερστερεοποιημένες (Lightly OC) & 1.5 -- 2.0 & $0.25 \sim 0.5$ & $0.5 \sim 0.7$ \\
    Ισχυρά υπερστερεοποιημένες (Heavily OC) & $> 4.0$ & $-0.5 \sim 0.0$ & $0.2 \sim 0.5$ \\
    \bottomrule
  \end{tabular}
\end{table}

\subsection{Επίλυση 3D Στερεοποίησης (Μέθοδος Πεπερασμένων Διαφορών)}

\subsubsection{Παραγωγή της 3D Διαφορικής Εξίσωσης}

Η διαφορική εξίσωση της στερεοποίησης προκύπτει από τον συνδυασμό τριών θεμελιωδών φυσικών νόμων:

\textbf{1. Αρχή Διατήρησης της Μάζας (Εξίσωση Συνέχειας)}\\
Σε ένα στοιχειώδες κυβικό εδαφικό στοιχείο (πλήρως κορεσμένο), ο ρυθμός με τον οποίο ο όγκος του νερού εξέρχεται από το στοιχείο (απόκλιση της ταχύτητας) πρέπει να ισούται με τον ρυθμό μείωσης του όγκου του (ρυθμός ογκομετρικής παραμόρφωσης, $\epsilon_v$):
\begin{equation}
  \frac{\partial v_x}{\partial x} + \frac{\partial v_y}{\partial y} + \frac{\partial v_z}{\partial z} = \frac{\partial \epsilon_v}{\partial t}
\end{equation}

\textbf{2. Νόμος του Darcy (Κίνηση του Νερού)}\\
Η ταχύτητα ροής του νερού στους πόρους είναι ανάλογη της υδραυλικής κλίσης. Η υδραυλική κλίση $i$ προκύπτει από την κλίση του υδραυλικού φορτίου $h = \frac{u}{\gamma_w} + z$. Επειδή η γεωμετρική κλίση (βαρύτητα) ακυρώνεται από την υδροστατική πίεση, η ροή προκαλείται αποκλειστικά από την υπερπίεση πόρων $u$:
\begin{equation}
  v_x = -k_x \frac{\partial h}{\partial x} = -\frac{k_x}{\gamma_w}\frac{\partial u}{\partial x}
\end{equation}
\begin{equation}
  v_y = -\frac{k_y}{\gamma_w}\frac{\partial u}{\partial y}, \qquad v_z = -\frac{k_z}{\gamma_w}\frac{\partial u}{\partial z}
\end{equation}
Αντικαθιστώντας τον νόμο του Darcy στην εξίσωση συνέχειας:
\begin{equation}
  -\left[\frac{k_x}{\gamma_w}\frac{\partial^2 u}{\partial x^2} + \frac{k_y}{\gamma_w}\frac{\partial^2 u}{\partial y^2} + \frac{k_z}{\gamma_w}\frac{\partial^2 u}{\partial z^2}\right] = \frac{\partial \epsilon_v}{\partial t}
\end{equation}

\textbf{3. Καταστατική Σχέση Εδάφους (Συμπιεστότητα)}\\
Η ογκομετρική παραμόρφωση $\epsilon_v$ συνδέεται με τη μεταβολή της ενεργού τάσης $\sigma'$ μέσω του συντελεστή ογκομετρικής συμπιεστότητας $m_v$:
\begin{equation}
  \partial\epsilon_v = m_v\,\partial\sigma'
\end{equation}
Βάσει της αρχής ενεργών τάσεων του Terzaghi ($\sigma = \sigma' + u$), και υποθέτοντας ότι η συνολική εξωτερική τάση παραμένει σταθερή στον χρόνο ($\partial\sigma = 0$), έχουμε $\partial\sigma' = -\partial u$. Συνεπώς:
\begin{equation}
  \frac{\partial \epsilon_v}{\partial t} = -m_v \frac{\partial u}{\partial t}
\end{equation}

\textbf{Τελική Εξίσωση Στερεοποίησης}\\
Αντικαθιστώντας το ρυθμό παραμόρφωσης στην εξίσωση, καταλήγουμε στην περίφημη 3D διαφορική εξίσωση (Terzaghi / Biot simplified):
\begin{equation}
  m_v \frac{\partial u}{\partial t} = \frac{k_x}{\gamma_w}\frac{\partial^2 u}{\partial x^2} + \frac{k_y}{\gamma_w}\frac{\partial^2 u}{\partial y^2} + \frac{k_z}{\gamma_w}\frac{\partial^2 u}{\partial z^2}
\end{equation}

Ορίζοντας τους συντελεστές στερεοποίησης $c_v = \frac{k_y}{\gamma_w m_v}$ και $c_h = \frac{k_x}{\gamma_w m_v} = \frac{k_z}{\gamma_w m_v}$, η εξίσωση παίρνει την τελική της μορφή:
\begin{equation}
  \frac{\partial u}{\partial t} = c_h \frac{\partial^2 u}{\partial x^2} + c_v \frac{\partial^2 u}{\partial y^2} + c_h \frac{\partial^2 u}{\partial z^2}
\end{equation}

Επειδή όμως τα $k$ και $m_v$ μεταβάλλονται από κόμβο σε κόμβο (διαφορετικές στρώσεις, διαφορετικό επίπεδο τάσεων), το πρόγραμμα επιλύει την εξίσωση στη \textbf{συντηρητική} της μορφή
\begin{equation}
  \gamma_w m_v \frac{\partial u}{\partial t} = \nabla\cdot\left(\mathbf{k}\,\nabla u\right)
\end{equation}
με ένα \textbf{Πεπλεγμένο Σχήμα στον Χρόνο (Backward Euler) και Κεντρικές Διαφορές στον Χώρο}.

\subsubsection{Σχήμα Πεπερασμένων Διαφορών}\label{sec:fd}

Η παράγωγος του χρόνου προσεγγίζεται με προς τα πίσω διαφορά (Backward difference): οι χωρικές παράγωγοι υπολογίζονται στο \textbf{νέο} χρονικό βήμα $n+1$ και όχι στο τρέχον.
\begin{equation}
  \frac{\partial u}{\partial t} \approx \frac{u^{n+1}_{i,j,k} - u^{n}_{i,j,k}}{\Delta t}
\end{equation}

Οι χωρικές παράγωγοι προσεγγίζονται με κεντρικές διαφορές, σχηματίζοντας ένα \textbf{7-point stencil (αστέρι 7 σημείων)}. Στη συντηρητική μορφή η ροή μεταξύ δύο γειτονικών κόμβων υπολογίζεται με τη διαπερατότητα της διεπιφάνειάς τους, π.χ. κατά $x$:
\begin{equation}\label{eq:flux}
  \frac{\partial}{\partial x}\left(k\frac{\partial u}{\partial x}\right) \approx \frac{\bar{k}_{i+\frac12}\left(u_{i+1}-u_i\right) - \bar{k}_{i-\frac12}\left(u_i-u_{i-1}\right)}{\Delta x^2},
  \qquad \bar{k}_{i+\frac12} = \frac{2\,k_i\,k_{i+1}}{k_i + k_{i+1}}
\end{equation}
όπου $\bar{k}$ ο \textbf{αρμονικός μέσος} των διαπερατοτήτων των δύο κόμβων. Για ομοιογενές έδαφος η σχέση ανάγεται στη γνωστή κεντρική διαφορά $\frac{u_{i-1} - 2u_i + u_{i+1}}{\Delta x^2}$.

\textbf{Αναπαράσταση του 3D Stencil:}
\begin{center}
\begin{tikzpicture}[x=1cm,y=1cm]
  \draw[thick] (-2.2,0) -- (2.2,0);
  \draw[thick] (0,-2.2) -- (0,2.2);
  \draw[dashed,gray] (-1.6,-1.6) -- (1.6,1.6);
  \filldraw[fill=blue!10,draw=blue!70!black,thick] (0,0) circle (0.45);
  \node at (0,0) {$u_{i,j,k}$};
  \node[above] at (0,2.2) {$u_{i,j-1,k}$ (Πάνω)};
  \node[below] at (0,-2.2) {$u_{i,j+1,k}$ (Κάτω)};
  \node[left] at (-2.2,0) {$u_{i-1,j,k}$ (Αριστερά)};
  \node[right] at (2.2,0) {$u_{i+1,j,k}$ (Δεξιά)};
  \node[gray,above right] at (1.5,1.5) {$u_{i,j,k-1}$ (Πίσω)};
  \node[gray,below left] at (-1.5,-1.5) {$u_{i,j,k+1}$ (Μπροστά)};
\end{tikzpicture}
\end{center}

Η αλγεβρική μορφή για κάθε κόμβο που δεν είναι κόμβος στράγγισης είναι:
\begin{equation}\label{eq:be}
  \gamma_w\, m_{v,i,j,k}\,\frac{u^{n+1}_{i,j,k} - u^{n}_{i,j,k}}{\Delta t} = \sum_{nb=1}^{6} \frac{\bar{k}_{nb}}{\Delta_{nb}^2}\left(u^{n+1}_{nb} - u^{n+1}_{i,j,k}\right)
\end{equation}
όπου το άθροισμα περιλαμβάνει τους 6 γείτονες του stencil, $\bar{k}_{nb}$ ο αρμονικός μέσος (με $k_h$ κατά $x$ και $z$, $k_v$ κατά $y$) και $\Delta_{nb}$ το αντίστοιχο βήμα του καννάβου. Επειδή οι άγνωστες τιμές $u^{n+1}$ εμφανίζονται και στα δύο μέλη, σε κάθε χρονικό βήμα προκύπτει ένα γραμμικό σύστημα
\begin{equation}
  \left(\frac{\mathbf{S}}{\Delta t} + \mathbf{K}\right)\mathbf{u}^{n+1} = \frac{\mathbf{S}}{\Delta t}\,\mathbf{u}^{n}
\end{equation}
όπου $\mathbf{S}$ το διαγώνιο μητρώο αποθηκευτικότητας ($\gamma_w m_v$) και $\mathbf{K}$ το μητρώο αγωγιμοτήτων. Το μητρώο του συστήματος είναι συμμετρικό και θετικά ορισμένο· το σύστημα λύνεται επαναληπτικά με τη μέθοδο των \textbf{Προρρυθμισμένων Συζυγών Κλίσεων (PCG)} με διαγώνιο προρρυθμιστή (Jacobi), έως σχετικό υπόλοιπο $10^{-7}$.

\subsubsection{Δυναμικός Συντελεστής Στερεοποίησης \& Διεπιφάνειες Στρώσεων}\label{sec:cv}

Ο συντελεστής $c_v$ υπολογίζεται σε \textbf{κάθε κόμβο} του καννάβου δυναμικά, βάσει της διαπερατότητας $k_v$ και του συντελεστή συμπιεστότητας $m_v$:
\begin{equation}
  c_v = \frac{k_v}{\gamma_w\, m_v}
\end{equation}

Ο $m_v$ υπολογίζεται σε κάθε κόμβο από την καμπύλη e-log(p') της στρώσης στην οποία ανήκει. Όπου υπάρχει πρόσθετη τάση ($\Delta\sigma_v > 1$ kPa) λαμβάνεται ως \textbf{τέμνον} μέτρο από την αρχική έως την τελική τάση, με το $\Delta e$ των εξισώσεων~(\ref{eq:reload})--(\ref{eq:mixed}):
\begin{equation}
  m_v = \frac{\Delta e}{(1+e_0)\,\Delta\sigma_v}
\end{equation}
Όπου η πρόσθετη τάση είναι αμελητέα, λαμβάνεται ως \textbf{εφαπτόμενο} μέτρο στην αρχική τάση, με $C = C_s$ για OCR $>1$ και $C = C_c$ για OCR $=1$:
\begin{equation}
  m_v = \frac{C}{\ln 10\,(1+e_0)\,\sigma'_{v0}}
\end{equation}
Ο συντελεστής $\mu$ (§\ref{sec:mu}) και η διόρθωση άκαμπτου θεμελίου (§\ref{sec:kany}) δεν μπαίνουν στον $m_v$: διορθώνουν το μέγεθος της καθίζησης και όχι τον ρυθμό της.

Αυτό επιτρέπει τη μοντελοποίηση ανομοιογενούς στερεοποίησης (π.χ. διαφορετικό $c_v$ κάτω από το κέντρο του φορτίου σε σχέση με τα άκρα, ανάλογα με το τοπικό επίπεδο τάσεων). Καθώς η $\sigma'_{v0}$ αυξάνει με το βάθος, ο $m_v$ μειώνεται και ο $c_v$ αυξάνει, οπότε τα βαθύτερα τμήματα μιας ομοιογενούς στρώσης στερεοποιούνται ταχύτερα.

\textbf{Αντιμετώπιση Διεπιφανειών με διαφορετικό $k$:} Σε περιπτώσεις όπου το μοντέλο περιλαμβάνει εναλλασσόμενες στρώσεις με διαφορετικές διαπερατότητες (π.χ. άργιλος με $k = 10^{-9}$ m/s πάνω από ιλύ με $k = 10^{-7}$ m/s), η ροή μεταξύ δύο γειτονικών κόμβων υπολογίζεται με τον \textbf{αρμονικό μέσο} των διαπερατοτήτων τους (εξ.~\ref{eq:flux}), που είναι η ισοδύναμη διαπερατότητα δύο στρώσεων σε σειρά. Έτσι η ροή που εξέρχεται από τον έναν κόμβο εισέρχεται ακριβώς στον γειτονικό (διατήρηση της μάζας) και η δυσπερατή στρώση ελέγχει σωστά τη ροή στη διεπιφάνεια. Συνεπώς, η πίεση διαχέεται γρήγορα στη στρώση με το υψηλό $k$, λειτουργώντας ως φυσική οδός αποστράγγισης για την παρακείμενη δυσπερατή στρώση (χαμηλό $k$), εξασφαλίζοντας αυτόματα τη συνέχεια της ροής στην υπολογιστική διεπιφάνεια.

\subsubsection{Τυπικές Τιμές Διαπερατότητας ($k_v$)}

Ο καθορισμός της ορθής τιμής κατακόρυφης διαπερατότητας ($k_v$) είναι κρίσιμος για τον υπολογισμό του $c_v$ και κατ' επέκταση για την ακριβή πρόβλεψη του χρόνου στερεοποίησης. Το λογισμικό δέχεται τιμές σε m/s. Στον Πίνακα~\ref{tab:k} παρουσιάζονται ενδεικτικές τυπικές τιμές για διάφορα είδη εδαφών.

\begin{table}[H]
  \centering
  \caption{Τυπικές τιμές κατακόρυφης διαπερατότητας ($k_v$) εδαφών}\label{tab:k}
  \begin{tabular}{lc}
    \toprule
    \textbf{Τύπος Εδάφους} & \textbf{Διαπερατότητα $k_v$ (m/s)} \\
    \midrule
    Καθαρά Χαλίκια (Gravels) & $> 10^{-2}$ \\
    Καθαρές Άμμοι (Sands) & $10^{-2} \sim 10^{-5}$ \\
    Ιλυώδεις Άμμοι / Ιλύες (Silts) & $10^{-5} \sim 10^{-9}$ \\
    Κανονικές Άργιλοι (Clays) & $10^{-9} \sim 10^{-11}$ \\
    Συμπαγείς Άργιλοι (Stiff Clays) & $< 10^{-11}$ \\
    \bottomrule
  \end{tabular}
\end{table}

\subsubsection{Οριακές Συνθήκες \& Βραχώδες Υπόβαθρο}\label{sec:bc}

\textbf{Α. Οριακές Συνθήκες Στράγγισης (Dirichlet - Prescribed Head):}
\begin{itemize}
  \item \textbf{Ελεύθερη Επιφάνεια:} Στους κόμβους που βρίσκονται στην επιφάνεια του εδάφους επιβάλλεται μόνιμα $u^{t}_{i,j_{surf},k} = 0$.
  \item \textbf{Ζώνη πάνω από τη Γ.Ο.Υ.:} Το έδαφος πάνω από τη στάθμη του υπόγειου νερού στραγγίζει ελεύθερα: στους κόμβους του επιβάλλεται $u = 0$, οπότε το όριο στράγγισης βρίσκεται στη Γ.Ο.Υ. Εκεί δεν αναπτύσσεται υπερπίεση και η καθίζηση θεωρείται άμεση ($U_z = 100\%$).
  \item \textbf{Στοιχεία Στράγγισης (Drains):} Σε όσους κόμβους απέχουν από μια γραμμή στράγγισης που σχεδιάζει ο χρήστης λιγότερο από τη μισή διαγώνιο του κελιού του καννάβου, επιβάλλεται $u = 0$. Έτσι κάθε γραμμή στράγγισης αντιστοιχεί σε μία συνεχή σειρά κόμβων.
\end{itemize}

\textbf{Β. Αδιαπέρατα Όρια (Neumann - Zero Flux):}\\
Στα πλευρικά όρια του μοντέλου (και στον πυθμένα, αν δεν υπάρχει στράγγιση), επιβάλλεται συνθήκη μηδενικής ροής $\frac{\partial u}{\partial n} = 0$. Αυτό υλοποιείται στον αλγόριθμο μέσω της μεθόδου των \textbf{φανταστικών κόμβων (ghost nodes)}. Για παράδειγμα, στο αριστερό όριο ($i = 0$), ο φανταστικός κόμβος $i = -1$ τίθεται ίσος με τον εσωτερικό κόμβο $i = 1$ ($u_{-1} = u_1$). Η κεντρική διαφορά γίνεται:
\begin{equation}
  \frac{u_1 - 2u_0 + u_{-1}}{\Delta x^2} \;\rightarrow\; \frac{2u_1 - 2u_0}{\Delta x^2}
\end{equation}
Ισοδύναμα, ο συνοριακός κόμβος αντιπροσωπεύει μισό όγκο ελέγχου (μισό κελί σε κάθε διεύθυνση όπου βρίσκεται σε όριο), ώστε το αδιαπέρατο όριο να βρίσκεται ακριβώς στον ακραίο κόμβο.

\textbf{Γ. Προσομοίωση Βραχώδους Υποβάθρου (Bedrock):}\\
Ανάλογα με το γεωλογικό προφίλ, η βάση του μοντέλου συμπεριφέρεται διαφορετικά:
\begin{itemize}
  \item \textbf{Αδιαπέρατο Υπόβαθρο:} Εάν η βάση του μοντέλου εδράζεται σε συμπαγή άρρηκτο βράχο ή πολύ σκληρή δυσπερατή άργιλο, \textbf{δεν} σχεδιάζουμε καμία γραμμή στράγγισης στον πυθμένα. Ο αλγόριθμος εφαρμόζει αυτόματα τη συνθήκη μηδενικής ροής (zero flux), επιβάλλοντας μονόπλευρη στερεοποίηση (τα ύδατα οδεύουν μόνο προς την επιφάνεια).
  \item \textbf{Διαπερατό Υπόβαθρο:} Εάν κάτω από τις αργιλικές στρώσεις υπάρχει υδροφόρος ορίζοντας, άμμος, χαλίκι ή ρηγματωμένος βράχος που δρα ως στραγγιστήρι, ο χρήστης πρέπει να σχεδιάσει μια \textbf{Γραμμή Στράγγισης} ακριβώς στη βάση του μοντέλου ($y = 0$). Το πρόγραμμα τότε θα επιβάλει $u = 0$ στους κάτω κόμβους, επιτρέποντας την αμφίπλευρη στερεοποίηση του υπερκείμενου συμπιεστού στρώματος.
\end{itemize}

\subsubsection{Ευστάθεια \& Χρονική Διακριτοποίηση}

Το πεπλεγμένο σχήμα Backward Euler είναι \textbf{ευσταθές για κάθε $\Delta t$}. Δεν υπάρχει κριτήριο τύπου Courant (CFL), όπως στα άμεσα σχήματα (FTCS), όπου θα έπρεπε
\begin{equation}
  \Delta t \le \frac{1}{2\left(\frac{c_h}{\Delta x^2} + \frac{c_v}{\Delta y^2} + \frac{c_h}{\Delta z^2}\right)_{max}}
\end{equation}
Αυτό έχει μεγάλη πρακτική σημασία: σε άμεσο σχήμα, μία μόνο διαπερατή στρώση (μεγάλο $c_v$) θα επέβαλλε εξαιρετικά μικρό $\Delta t$ σε όλο το μοντέλο. Εδώ το βήμα του χρόνου δεν εξαρτάται από τον $c_v$, οπότε οι διαπερατές στρώσεις προσομοιώνονται με την πραγματική τους διαπερατότητα, χωρίς τεχνητό περιορισμό.

Τα χρονικά βήματα είναι \textbf{λογαριθμικά ισαπέχοντα} (30 ανά δεκάδα χρόνου, από 1 ώρα έως τον μέγιστο χρόνο), ώστε να αποδίδεται με ακρίβεια τόσο η απότομη εκτόνωση των πρώτων ημερών όσο και η αργή εξέλιξη των τελευταίων ετών. Στα βήματα περιλαμβάνονται ακριβώς οι χρόνοι των \textbf{40 στιγμιοτύπων} της λύσης, που ισαπέχουν λογαριθμικά από 1 ημέρα έως τον μέγιστο χρόνο.

\subsubsection{Βαθμός Στερεοποίησης \& Καθίζηση στον Χρόνο}

Ο τοπικός βαθμός στερεοποίησης σε κάθε κόμβο ορίζεται από την αρχική υπερπίεση $u_0$ ($= \Delta\sigma_v$ κάτω από τη Γ.Ο.Υ., μηδέν πάνω από αυτήν):
\begin{equation}
  U_z(t) = 1 - \frac{u(t)}{u_0}
\end{equation}
Όπου $u_0 = 0$ η καθίζηση είναι άμεση ($U_z = 100\%$). Η καθίζηση στον χρόνο $t$ προκύπτει αθροίζοντας τις συμπιέσεις $\Delta S_j$ των διαστημάτων μεταξύ διαδοχικών κόμβων κατά το βάθος, πολλαπλασιασμένες με τον μέσο βαθμό στερεοποίησης των δύο άκρων κάθε διαστήματος (τραπεζοειδής κανόνας):
\begin{equation}
  S(t) = \sum_{j} \Delta S_j\,\frac{U_{z,j}(t) + U_{z,j+1}(t)}{2}
\end{equation}

Ο κόμβος στράγγισης στερεοποιείται ακαριαία, άρα ο μισός όγκος ελέγχου γύρω του (πάχους $\Delta y/2$) εμφανίζεται πλήρως στερεοποιημένος από την πρώτη στιγμή. Κοντά στην επιφάνεια, όπου η $\sigma'_{v0}$ είναι μικρή και οι παραμορφώσεις μεγάλες, αυτό οδηγεί σε υπερεκτίμηση της καθίζησης των πρώτων ημερών· η υπερεκτίμηση μειώνεται όσο πυκνώνει ο κάνναβος κατά το βάθος, ενώ η καθίζηση στους μεγαλύτερους χρόνους δεν επηρεάζεται.

\subsubsection{Επαλήθευση έναντι της Λύσης Terzaghi}

Ο επιλυτής ελέγχθηκε σε μονοδιάστατη στήλη ύψους $H = 20$ m (41 κόμβοι, όπως στην εφαρμογή) με ομοιογενές έδαφος, στραγγιζόμενη επιφάνεια και είτε αδιαπέρατη είτε στραγγιζόμενη βάση. Ο μέσος βαθμός στερεοποίησης συγκρίνεται με την αναλυτική λύση του Terzaghi στον Πίνακα~\ref{tab:terzaghi}. Η απόκλιση δεν υπερβαίνει το $0.7\%$ και οφείλεται κυρίως στη χρονική διακριτοποίηση πρώτης τάξης.

\begin{table}[H]
  \centering
  \caption{Μέσος βαθμός στερεοποίησης $U$ (\%): \cons{} έναντι αναλυτικής λύσης Terzaghi}\label{tab:terzaghi}
  \begin{tabular}{cccc}
    \toprule
    \textbf{$T_v$} & \textbf{Terzaghi} & \textbf{Μονόπλευρη στράγγιση} & \textbf{Αμφίπλευρη στράγγιση} \\
    \midrule
    0.01 & 11.28 & 11.28 & 11.41 \\
    0.05 & 25.23 & 25.16 & 25.21 \\
    0.10 & 35.68 & 35.57 & 35.61 \\
    0.20 & 50.41 & 50.22 & 50.24 \\
    0.50 & 76.40 & 75.90 & 75.90 \\
    1.00 & 93.13 & 92.51 & 92.50 \\
    \bottomrule
  \end{tabular}
\end{table}

% =====================================================================
\section{Μέρος 2: Εγχειρίδιο Χρήσης}

\subsection{Διεπαφή \& Εργαλεία (Toolbar)}

Στην αριστερή πλευρά θα βρείτε τα εργαλεία σχεδίασης:
\begin{itemize}
  \item \textbf{Επιλογή (MousePointer)}: Για επιλογή υπαρχόντων στοιχείων και επεξεργασία των ιδιοτήτων τους στο δεξί πάνελ.
  \item \textbf{Τομή X-Y}:
  \begin{itemize}
    \item \textbf{Στρώση (Square)}: Σχεδίαση πολυγωνικών στρώσεων εδάφους.
    \item \textbf{2D Επίχωμα (Pencil)}: Ορισμός επιχώματος που εκτείνεται σε όλο το μήκος Z (λωριδωτό φορτίο).
    \item \textbf{Στράγγιση (Waves)}: Σχεδίαση γραμμών που λειτουργούν ως ανοιχτά σύνορα (\mbox{$u = 0$}).
  \end{itemize}
  Η σχεδίαση ενός πολυγώνου ή μιας γραμμής ολοκληρώνεται με δεξί κλικ.
  \item \textbf{Κάτοψη X-Z}:
  \begin{itemize}
    \item \textbf{Ορθογώνιο/Κυκλικό Φορτίο}: Τοποθέτηση 3D φορτίων στην επιφάνεια.
    \item \textbf{3D Επίχωμα (Scan)}: Δημιουργία επιχώματος με συγκεκριμένο μήκος και πλάτος (πυραμιδοειδές).
  \end{itemize}
\end{itemize}

\subsection{Ρύθμιση Μοντέλου \& Τρισδιάστατη (3D) Απεικόνιση}

Στο δεξί πάνελ της εφαρμογής βρίσκονται οι κεντρικές ρυθμίσεις του μοντέλου, αλλά και οι επιλογές της ισχυρής μηχανής 3D απεικόνισης.
\begin{itemize}
  \item \textbf{Στάθμη Γ.Ο.Υ}: Καθορίζει την άνωση και τις αρχικές ενεργές τάσεις $\sigma'_{v0}$, καθώς και το άνω όριο στράγγισης της στερεοποίησης (§\ref{sec:bc}).
  \item \textbf{Μήκος Z}: Ορίζει το βάθος της 3ης διάστασης για τον υπολογισμό των 3D επιρροών και τη διακριτοποίηση.
  \item \textbf{Παραμορφωμένο Δίκτυο}: Ενεργοποιεί την οπτική αναπαράσταση της βύθισης του εδάφους (με slider μεγέθυνσης για εποπτεία).
  \item \textbf{Τρισδιάστατη (3D) Απεικόνιση Ισοκαμπυλών (Heatmaps):} Ένα από τα σημαντικότερα και πιο καινοτόμα χαρακτηριστικά του \cons{} είναι η δυνατότητα δυναμικής 3D απεικόνισης των αποτελεσμάτων. Το λογισμικό επιτρέπει την προβολή ισοκαμπυλών για κάθε υπολογιζόμενο μέγεθος ($\Delta\sigma_v, S, u, U_z$ κ.ά.) ταυτόχρονα σε \textbf{τρία κάθετα επίπεδα} (X-Y, X-Z, Y-Z).

  Το κρίσιμο πλεονέκτημα είναι ότι \textbf{η θέση αυτών των επιπέδων είναι ελεύθερα μεταβαλλόμενη}. Μέσω ειδικών ρυθμιστών (sliders), ο χρήστης μπορεί να ``σαρώσει'' δυναμικά όλο τον όγκο του εδάφους, μετακινώντας τα επίπεδα προβολής κατά τον άξονα X, Y ή Z. Αυτό προσφέρει την απόλυτη χωρική εποπτεία της κατανομής των τάσεων (π.χ. πώς αναπτύσσεται ο βολβός τάσεων κάτω από ένα 3D επίχωμα) και της εξέλιξης των υπερπιέσεων πόρων στο χώρο και τον χρόνο.
\end{itemize}

\subsection{Εκτέλεση Ανάλυσης Χρόνου}

\begin{enumerate}
  \item Ορίστε την κατακόρυφη διαπερατότητα ($k_v$) και την ανισοτροπία ($k_h/k_v$) στις ιδιότητες των στρώσεων εδάφους.
  \item Στο πεδίο \textbf{Εξέλιξη στον Χρόνο}, ορίστε τον μέγιστο χρόνο προσομοίωσης σε έτη (δεκτές και δεκαδικές τιμές, π.χ. 0.5).
  \item Πατήστε \textbf{Επίλυση}. Το πεπλεγμένο σχήμα (§\ref{sec:fd}) εκτελεί τη χρονική ολοκλήρωση και αποθηκεύει 40 στιγμιότυπα της λύσης· η επίλυση διαρκεί συνήθως λίγα δευτερόλεπτα.
  \item Μετά την επίλυση, χρησιμοποιήστε τον ρυθμιστή \textbf{Timeline Στερεοποίησης} για να δείτε πώς μεταβάλλονται οι πιέσεις πόρων ($\Delta u$) και οι καθιζήσεις από την 1η ημέρα έως τον μέγιστο χρόνο. Η τελευταία θέση του ρυθμιστή αντιστοιχεί σε άπειρο χρόνο ($\Delta u = 0$, τελική καθίζηση).
\end{enumerate}

\subsection{Εργαλεία Λεπτομερούς Ανάλυσης}

\begin{itemize}
  \item \textbf{e-Log(p') Analysis (Calculator)}: Κάντε κλικ σε οποιοδήποτε σημείο της τομής (X-Y) για να ανοίξει το αναδυόμενο παράθυρο. Εκεί εμφανίζεται οπτικοποιημένη η αναλυτική πορεία συμπίεσης, ο υπολογισμός του $p'_c$ και η μετάβαση από τον κλάδο ανασυμπίεσης στην παρθένα καμπύλη.
  \item \textbf{Time Plot (LineChart)}: Κάντε κλικ στο πλέγμα για να εξαγάγετε το διάγραμμα εξέλιξης της επιφανειακής καθίζησης και της εσωτερικής υπερπίεσης πόρων ($u$) στο συγκεκριμένο κατακόρυφο προφίλ. Ο άξονας του χρόνου είναι λογαριθμικός, από 0.1 ημέρα (αρχική κατάσταση) έως την επόμενη ακέραια δεκάδα ημερών· με το ποντίκι πάνω στο διάγραμμα εμφανίζονται οι τιμές κάθε στιγμιοτύπου.
\end{itemize}

% =====================================================================
\section{Μέρος 3: Συμβουλές για Ακριβή Μοντελοποίηση}

\begin{enumerate}
  \item \textbf{Άκαμπτες έναντι Εύκαμπτων Θεμελιώσεων}: Όταν θέλετε να προσομοιώσετε κοιτόστρωση σκυροδέματος, ενεργοποιήστε το ``Άκαμπτη Θεμελίωση (Rigid)''. Η προσέγγιση του Χαρακτηριστικού Σημείου του Kany εξισορροπεί με ρεαλιστικό τρόπο το ``μπολ'' της εύκαμπτης καθίζησης.
  \item \textbf{Συντελεστής Ανισοτροπίας ($k_h/k_v$)}: Λόγω της εναπόθεσης των αργίλων στη φύση, συνήθως $k_h > k_v$ (συχνά $k_h \approx 2k_v$ έως $5k_v$). Η εισαγωγή αυτής της παραμέτρου επιταχύνει ρεαλιστικά τη στερεοποίηση, ειδικά κάτω από τοπικά (όχι απειρομήκη) φορτία όπου υπάρχει έντονη οριζόντια αποστράγγιση.
  \item \textbf{Οριακές Συνθήκες Στράγγισης}: Βεβαιωθείτε ότι αντιλαμβάνεστε πλήρως το υδρογεωλογικό καθεστώς. Η παράλειψη της τοποθέτησης γραμμής στράγγισης σε επαφή με ένα διαπερατό στρώμα/bedrock, θα αναγκάσει το λογισμικό να υποθέσει αδιαπέρατο όριο (zero flux) με αποτέλεσμα οι χρόνοι στερεοποίησης (τουλάχιστον στον κατακόρυφο άξονα) να τετραπλασιαστούν!
\end{enumerate}

\end{document}
